\section{引言}

硒是哺乳动物必需的微量元素，参与谷胱甘肽过氧化物酶等硒蛋白的合成，在维持机体氧化还原稳态中发挥重要作用。与亚硒酸盐、硒酸盐等无机硒相比，零价态纳米硒在部分研究中表现出较高的生物利用度和较低的急性毒性，在抗肿瘤、抗菌、抗糖尿病及抗氧化等生物医学领域具有应用潜力\cite{hosnedlova2018nanoselenium,bisht2022senps,sowmya2024nanoselenium,karthik2024senps}。然而，硒的有效剂量与毒性剂量较为接近，纳米化也不能完全消除其毒性风险。研究发现，纳米硒在发挥生物学作用的同时，仍可能引起氧化损伤和细胞毒性\cite{wang2020prooxidation}。过量硒暴露还可能造成消化系统、呼吸系统及心血管系统等损伤\cite{hosnedlova2018nanoselenium}。因此，预测纳米硒在不同材料性质和暴露条件下的毒性，是其生物医学应用的重要基础。

纳米硒的细胞毒性与粒径、暴露浓度、暴露时间及表面修饰等因素密切相关。粒径减小可增加颗粒的比表面积，影响其生物活性和细胞摄取过程，也可能会增加毒性风险\cite{zhang2023senpsbio}。在暴露浓度方面，Wang 等\cite{wang2020prooxidation}发现，大鼠肝细胞暴露于 24~$\mu$M 和 48~$\mu$M 纳米硒 24~h 后，细胞存活率分别降至 87.76\% 和 51.37\%，表现出剂量依赖性。表面修饰也会改变纳米硒与细胞的相互作用。例如，螺旋藻多糖（SPS）、江蓠多糖（GLP）及黄芪多糖（AASP）修饰的纳米硒在相关研究中表现出对肿瘤细胞的选择性杀伤作用，同时对正常细胞的毒性较低\cite{zhang2023senpsbio,ji2023astragalus}；季铵化壳聚糖--没食子酸修饰的纳米硒则在 Colo357 和 HaCaT 细胞中表现出较高的细胞毒性\cite{lunkov2023chitosan}。这些结果说明，不同纳米硒材料的毒性不能仅由单一性质判断，需要结合材料特征、细胞类型和暴露条件进行考虑。

传统纳米毒性评价主要依赖体外细胞实验和体内动物实验，但实验周期、成本及检测通量等因素限制了其在大量材料筛选中的应用\cite{ahmadi2024toxicity}。此外，纳米材料的粒径、形貌、表面性质和化学组成等因素共同影响毒性，其与生物效应之间的关系较为复杂\cite{ji2025rosml}。随着计算毒理学的发展，纳米定量构效关系（nano-QSAR）和定量结构--毒性关系（nano-QSTR）等方法被用于分析材料性质与毒性之间的关系，成为实验评价的补充\cite{roy2023tio2}。其中，随机森林、支持向量机和神经网络等机器学习方法已经被用于金属及金属氧化物纳米颗粒的细胞毒性预测\cite{zhang2023qsar}。另外，为了改善模型的预测表现，研究人员针对机器学习方法结合特征筛选、聚类分析和参数优化进行了进一步的研究\cite{sang2022mlcytotoxicity}。

已有研究展示了机器学习在纳米毒性预测中的应用潜力。Shirokii 等通过相关性特征筛选和网格搜索优化模型参数，构建的无机纳米材料细胞毒性预测模型在 10 折交叉验证中表现良好，$Q^2$ 达到 0.86\cite{shirokii2023quantitative}。Hosni 等的研究表明，LightGBM 在所比较的模型中具有较好的纳米颗粒细胞毒性分类性能\cite{hosni2025mlnp}。此外，Boruta-RF-ARM 组合方法识别了浓度、核壳材料及细胞类型等与毒性相关的重要因素\cite{yildirim2025rheumatoid}。然而，现有模型多面向多种纳米材料，对纳米硒及其表面修饰差异的考虑仍不充分。因此，有必要结合纳米硒的材料性质与暴露条件建立细胞毒性预测模型，为实验前的毒性评估和材料初步筛选提供参考。

基于上述问题，本研究以细胞存活率为预测目标，进行以下研究贡献：
\begin{itemize}
\item[1)]利用通用纳米颗粒和纳米硒文献数据，比较不同机器学习模型的预测性能，结合材料化学特征、特征筛选和多模态信息开展建模;
\item[2)]利按文献来源分组评价模型的泛化能力，考察迁移学习及浓度与暴露时间交互项对预测性能的影响;
\item[3)]利开展 9 种纳米硒材料的人脐静脉内皮细胞（HUVEC）实验，通过预测结果与实测结果的比较，检验模型对新材料细胞存活率及毒性差异的预测能力。
\end{itemize}

本研究旨在探索依据材料性质和暴露条件预测纳米硒细胞毒性的方法，为实验前的毒性评估与材料初步筛选提供参考。
